\chapter{Mensageria}

A mensageria é o mecanismo que viabiliza o desacoplamento entre o início da análise e o processamento computacional subsequente. O microserviço utiliza RabbitMQ como broker e MassTransit como abstração de alto nível para simplificar publicação, consumo, configuração e evolução futura da topologia de filas \cite{rabbitmqdocs,masstransitdocs}.

\section{RabbitMQ e MassTransit}

RabbitMQ fornece o meio de transporte assíncrono, enquanto MassTransit simplifica a integração ao encapsular detalhes de conexão, roteamento e endpoints. Essa combinação é adequada para um cenário em que o sistema precisa iniciar análises sem bloquear a resposta HTTP e, ao mesmo tempo, preservar rastreabilidade e resiliência.

\section{Commands e Events}

O contrato de solicitação utiliza \textit{AnalysisRequestedCommand}, que representa a intenção de iniciar uma análise. Após o processamento, o sistema pode gerar \textit{AnalysisCompletedEvent} ou \textit{AnalysisFailedEvent}, refletindo os estados finais do ciclo de vida da análise. Essa separação semântica evita ambiguidade e melhora a compreensão da topologia de mensagens.

\section{Consumer e Processor}

O \textit{AnalysisRequestedConsumer} atua como adaptador entre a fila e a aplicação. Ele não contém regra de negócio e apenas encaminha o identificador da análise para o \textit{AnalysisProcessor}. O processor, por sua vez, concentra a orquestração: consulta o repositório, atualiza status, aciona o pipeline, salva o resultado e solicita o callback.

\section{Fluxo assíncrono}

O fluxo assíncrono reduz o acoplamento entre a requisição inicial e o processamento pesado. A API pode responder rapidamente, enquanto a execução do pipeline ocorre fora do caminho síncrono. Esse modelo melhora a escalabilidade e a experiência de integração com sistemas externos.

\section{Callback}

O callback é o mecanismo pelo qual o sistema solicitante recebe o resultado final da análise. Em vez de exigir polling contínuo, o microserviço pode notificar o backend originador quando a análise é concluída ou falha. Isso reduz o tráfego desnecessário e melhora a ergonomia de integração.

\begin{figure}[H]
    \centering
    \fbox{\parbox{0.85\textwidth}{\centering Espaço reservado para o diagrama de sequência da mensageria.}}
    \caption{Sequência assíncrona da análise}
    \label{fig:sequencia-mensageria}
\end{figure}
